\textit{Key Findings.} In this work, we analyzed interacting language-model agents over $10,000$ simulated groups and found structured collective dynamics across models, tasks, and communication networks. Repeated interaction increases conviction and shifts initially indifferent groups toward more ordered states. On objective questions, we find that collective accuracy improves over rounds. Meanwhile, on subjective questions, three of the four models exhibit a rightward drift on the political opinions. We then develop a statistical model of the mechanics of opinion updates. Our model, which we fit on a set of training questions, generalizes to unseen questions and graphs, predicts individual trajectories, and approximately reproduces group-level outcomes. The fitted parameters suggest that (1) the groups operate below a critical social temperature, which drives conviction buildup; (2) concordant interactions are stronger than the discordant interactions, which drives consensus formation; and (3) greater influence from correct neighbors helps explain truth-seeking on objective tasks. Together, these findings show that collective agent behavior can be both predictable and interpretable, while also highlighting that the effects of interaction depend on the setting because the same mechanisms can improve factual decisions or amplify political biases.

\textit{Limitations.} Our setup has three main simplifying assumptions: (1) The community answers a single shared binary question at a time, and each agent's output is summarized by a binary opinion state $s_i(t) \in \{+1,-1\}$ together with a short natural-language message. (2) The communication pattern is fixed across timesteps and symmetric, so who talks to whom does not change during a run. (3) An agent's update depends only on its persona, the question, and the messages in its current inbox, not the full history of interactions it has had. The main limitation of our method is that it simplifies the nature of the interactions by discarding the content of the natural-language messages that mediate influence.

\textit{Future Work (Setup).} These limitations point to natural extensions of our setup. (1) Opinions can be multiple choice answers, Likert scale \citep{likert1932}, or open ended statements in natural language, all of which can be represented in a vector space. (2) The agents in a community can discuss multiple questions and/or statements simultaneously, and one can study how the conversation changes opinions about multiple statements in conjunction. (3) The communication pattern can change at each timestep, or be even more flexible, allowing the agents to choose whom to communicate with in the next step. It can also be interesting to explore an extension in which agents are situated in a physical space and communicate only with others in close proximity. (4) Finally, memory and continual learning mechanisms can allow the agents to remember and learn from the experiences from the previous timesteps. (5) A similar analysis can be applied to data generated by AI agents in the wild on social platforms such as \href{https://www.moltbook.com/}{Moltbook} \citep{moltbook_observatory_archive_2026}.(6) Future work could also study larger and heterogeneous groups of language model agents, including communities composed of different language models, to understand how collective dynamics and finite-size effects change with population size and model diversity.
 
\textit{Future Work (Method).}  These extensions of the setup motivate corresponding extensions of our method. (1) A Potts-type generalization \citep{potts1952, wu1982potts} that treats opinions as vectors would extend the model from a single binary question to questions with open-ended or multiple candidate answers. (2) Message content can be incorporated into the local field, e.g., through message embeddings in the spirit of message passing in graph neural networks \citep{scarselli2009gnn, gilmer2017mpnn}, so that the couplings act on what a neighbor says rather than only on the sign of its stance. (3)  Beyond fitting observed dynamics, controlled interventions on edges, messages, or initial opinions could test whether the learned parameters predict how collective outcomes change under perturbations and could ultimately help design communication structures that promote desirable behavior \citep{hu2025automateddesignagenticsystems}.  To that end, learning rules, such as Hebbian Learning \citep{hebb1949organization}, can allow the communication patterns to evolve during the interactions.

\textit{Social Impacts.} Our work is a step toward understanding the behavior of multi-agent systems without having to run them at scale. First, a theory of collective dynamics can offer a way to anticipate failure modes in a world where AI agents routinely interact with one another, which is an increasingly likely scenario as AI agents become more widespread. Second, the same theory can inform the design of agentic systems. If collective outcomes are predictable, then communication channels, population composition, and other hyperparameters can become design variables that can be optimized theoretically. Finally, we caution against over-extrapolating from simulated persona-conditioned agents to human communities. Our agents are language models prompted to hold personas or expertise. The dynamics we measure characterize that system, and any resemblance to human opinion dynamics is a hypothesis for future work.
